\chapter{Derived dg categories and Morita theory}
\label{ch:dg-categories}

The representable resolutions above let us compute the derived
category without any boundedness condition. Its compact objects
recover the finite shifts, cones, and retracts of the representables.
We then use the signed tensor--Hom adjunction to compare derived
categories of different dg categories.

\section{Homotopy-projectivity and the derived category}

\begin{theorem}[Semi-free modules are \(K\)-projective]\label{res:G07}
A dg module \(P\) is called \(K\)-projective if
\(\underline{\operatorname{Hom}}_{\mathcal A}(P,N)\) is acyclic for
every acyclic \(N\). Every semi-free module is \(K\)-projective.
The class of \(K\)-projective modules is closed under shifts,
cones of closed degree-zero maps, and arbitrary direct sums.
\end{theorem}
\begin{proof}
Let \(N\) be acyclic and let \(f:P\to N\) be a graded map of
degree \(r\) with \(\partial f=0\), where \(P\) is semi-free.
Construct a graded map \(h\) of degree \(r-1\) on successive
cell stages with \(\partial h=f\). At a new generator \(e\)
of degree \(n\) at \(a\), take \(z=de\) in the previous stage.
The requirement is
\[
          d_N h(e)= f(e)-(-1)^r h(z).
\]
The right side is closed because
\(d_N f(e)=(-1)^r f(z)\), and the already constructed homotopy
gives \(d_N h(z)=f(z)\) since \(dz=0\). Acyclicity of \(N(a)\)
gives a primitive, which we choose as \(h(e)\). Extending
right-linearly over the new cells is well-defined, and verification
on generators suffices by Result~\ref{res:G03}. Taking the union yields
\(\partial h=f\) at every degree. This proves \(K\)-projectivity.

Shifts preserve the condition by reindexing Hom complexes with
the shift signs of Result~\ref{res:G04}. For a cone \(C\) of
\(P_1\to P_2\), the degreewise split sequence
\(P_2\to C\to P_1[1]\) gives a short exact sequence of Hom complexes
into \(N\), in the reverse order, since its splitting exists
in graded \(\mathcal A\)-modules. The end complexes are acyclic.
For a sum of \(P_i\), Hom becomes the product of their Hom
complexes. A product of acyclic complexes of modules is acyclic
by choosing a primitive in every component of a cycle.
\end{proof}

\begin{theorem}[Computing the derived category by resolutions]\label{res:G08}
Let \(D(\mathcal A)\) be the localization of \(K(\mathcal A)\)
at the objectwise quasi-isomorphisms. The full homotopy subcategory
of \(K\)-projective modules maps equivalently onto \(D(\mathcal A)\).
If \(P\to M\) is a \(K\)-projective replacement, then
\[
 \operatorname{Hom}_{D(\mathcal A)}(M,N)
      \cong H^0\underline{\operatorname{Hom}}_{\mathcal A}(P,N).
\]
Cones of maps of \(K\)-projective modules represent distinguished
triangles, and coproducts of arbitrary derived objects are
represented by direct sums of their modules.
\end{theorem}
\begin{proof}
For a quasi-isomorphism \(u:N'\to N\) and \(K\)-projective \(P\),
\(\underline{\operatorname{Hom}}_{\mathcal A}(P,u)\) is a
quasi-isomorphism. Its cone is identified with Hom into the cone
of \(u\) using the Hom and shift formulas, hence is acyclic.
An acyclic \(K\)-projective module is contractible since its
identity represents a class in \(H^0\) of its own acyclic Hom
complex. By the closure in Result~\ref{res:G07}, the cone of a
quasi-isomorphism between \(K\)-projective modules is contractible,
so that quasi-isomorphism is a homotopy equivalence.

For every \(M\) choose \(p_M:P_M\to M\) as in Result~\ref{res:G06},
taking the identity when \(M\) is already \(K\)-projective.
For a homotopy class \(M\to N\), the isomorphism
\(\operatorname{Hom}_{K(\mathcal A)}(P_M,P_N)\to
  \operatorname{Hom}_{K(\mathcal A)}(P_M,N)\)
induced by \(p_N\) gives a unique lifted homotopy class
\(P_M\to P_N\). Uniqueness proves the composition law and
naturality of the \(p_M\). If the original map is a quasi-isomorphism,
so is its lift by the comparison square, and this lift is
invertible in the homotopy category. The lift functor therefore
factors through the localization. The comparisons \(p_M\) give a
natural isomorphism from its composite with inclusion to the
identity on the localization. In the other order the composite
is the identity on the \(K\)-projective subcategory. This proves
the equivalence and the stated Hom formula, since \(N\) can be
compared with \(P_N\) after applying Hom from \(P_M\).

The \(K\)-projective homotopy subcategory is triangulated by
Results~\ref{res:G04} and~\ref{res:G07}. It gives the cone
triangulation on the localization. Indeed a closed map between
arbitrary modules is represented by the lifted map between
replacements. A comparison with its original cone exists from
the degree-\(-1\) homotopy in the comparison square and the
usual cone matrix in \cite[K05]{ha-core}. It is a quasi-isomorphism
by the long exact sequences for the two cones. Direct sums of
\(K\)-projective modules are \(K\)-projective, and direct sums preserve
objectwise quasi-isomorphisms. Hom from a sum of replacements is
the product of the corresponding Hom groups, as \(H^0\) commutes
with products of complexes of modules. Hence the sum of the
chosen replacements computes the coproduct, and its map to the
sum of original modules is a quasi-isomorphism.
\end{proof}

\section{Compactness and the pretriangulated hull}

\begin{lemma}[The telescope of a sequence of strict maps]\label{res:G09}
Let \(M_0\to M_1\to\cdots\) be a sequence of dg modules and closed
degree-zero maps. On \(U=\bigoplus_{j\geq0}M_j\) let \(s\) send
the \(j\)-th summand to the \((j+1)\)-st by the transition.
Then \(\operatorname{Cone}(1-s)\to\varinjlim_j M_j\) is an
objectwise quasi-isomorphism. If \(C\) is compact in
\(D(\mathcal A)\), any map from \(C\) to this telescope in
\(D(\mathcal A)\) comes from a single \(M_j\). In fact its group
of maps is the colimit of the groups of maps into \(M_j\).
\end{lemma}
\begin{proof}
The map \(1-s\) is injective in each degree by considering the
first nonzero component of a finite-support element. Its cokernel
imposes exactly the relations identifying an element and its
successor, which presents the colimit. For any injective dg map
\(U'\to U''\), the natural map from its cone to its cokernel
has kernel the cone of the identity of \(U'\). That cone is
contractible by the right-linear homotopy
\((y,x)\mapsto(0,y)\). This proves the claim for \(1-s\).

Compactness means that \(\operatorname{Hom}(C,-)\) takes arbitrary
sums to direct sums, and this also applies to the objects \(M_j[1]\).
In the Hom exact sequence of the cone of \(1-s\), the map
\(1-s\) on the direct sum of the groups
\(\operatorname{Hom}(C,M_j[1])\) is injective by the same first
component argument. Thus \(\operatorname{Hom}(C,\operatorname{Cone}(1-s))\)
is the cokernel of \(1-s\) on the unshifted Hom groups. This is
their colimit. An element in a sequential colimit of groups
can be represented at a single stage.
\end{proof}

\begin{theorem}[Representables are compact generators]\label{res:G10}
Each \(h_a\) is a compact object of \(D(\mathcal A)\). For every
\(N\) and \(n\in\mathbb Z\),
\[
       \operatorname{Hom}_{D(\mathcal A)}(h_a[-n],N)
           \cong H^n(N(a)).
\]
The smallest full triangulated subcategory closed under
isomorphisms and coproducts and containing the \(h_a\) is all of
\(D(\mathcal A)\). In particular the shifted representables detect
the zero object.
\end{theorem}
\begin{proof}
The Hom formula follows from enriched Yoneda and
Results~\ref{res:G05} and~\ref{res:G08}. Explicitly a map out of
the shifted cell is determined by a degree-\(n\) element, closed
for a chain map, and homotopy identifies classes modulo boundaries.
Since cohomology of modules commutes with direct sums, the formula
also proves compactness by the coproduct description of
Result~\ref{res:G08}.

For generation resolve \(N\) by a semi-free module \(P\). Each
\(P_j/P_{j-1}\) is a direct sum of shifted representables, and
\(P_{j-1}\to P_j\to P_j/P_{j-1}\) is part of a triangle by the
cone-to-cokernel argument of Result~\ref{res:G09}. Therefore a
subcategory with the stated closure contains each \(P_j\).
It contains the telescope, which represents \(P\) by
Result~\ref{res:G09}, and hence contains \(N\). Alternatively,
vanishing of the displayed Hom groups says that \(N\) is
objectwise acyclic, so \(N\) is zero in the localization.
\end{proof}

\begin{theorem}[Compact objects and finite cells]\label{res:G11}
An object of \(D(\mathcal A)\) is compact if and only if it is
a retract in \(D(\mathcal A)\) of a semi-free module with only
finitely many cells.
\end{theorem}
\begin{proof}
Finite-cell objects are obtained by finite sums, shifts, and
cones from representables. Compact objects are closed under
those operations by the exact Hom sequences of triangles and
exactness of direct sums of groups. A retract is compact because
the coproduct comparison map for a retract is a retract of the
same map for the original object.

Conversely let \(C\) be compact and choose a semi-free
representative \(P\) of \(C\). By Result~\ref{res:G09} its identity
factors in \(D(\mathcal A)\) through one finite stage \(P_j\) of
the cell filtration, which may itself contain infinitely many
cells. At a fixed stage \(j\), call a finite subset of cells
closed if it contains all the cells occurring in the differential of
every cell in it. Every finite subset can be enlarged to a finite
closed subset by recursively adding the finite differential
supports, which lie at strictly earlier stages. Finite closed
subsets form a filtered system under unions. For such a subset
\(E\), let \(P_{j,E}\) be its finite-cell submodule.

For every \(t\) there is an isomorphism
\[
 \varinjlim_E\operatorname{Hom}_{D(\mathcal A)}(C,P_{j,E}[t])
     \ \cong\ \operatorname{Hom}_{D(\mathcal A)}(C,P_j[t]).
\]
Prove it by induction on \(j\). For stage zero it is compactness
on a sum of shifted representables, since each element of the sum
has finite support. At stage \(j\), the first \(j-1\) stages for a
finite \(E\) form \(P_{j-1,E\cap P_{j-1}}\), and the quotient is a
finite sum of cells from stage \(j\). The associated triangles,
including all shifts, give a comparison of long exact Hom sequences
after taking the filtered colimit. Filtered colimits of groups
are exact, since a finite relation between representatives holds
at a sufficiently late common stage. The group maps for the
lower-stage terms are isomorphisms by induction. The system of
intersections is cofinal among finite closed subsets of the
lower stage, since any such subset can itself be used as \(E\).
The quotient term maps are isomorphisms by compactness, since
each finite subset of top-stage cells can be enlarged to a finite
closed subset. Exactness on the two sides implies the middle
group maps are isomorphisms by the five-term diagram argument.
This proves the displayed comparison.

Consequently the factorization of the identity through \(P_j\)
factors further through a finite-cell \(P_{j,E}\). Its composite
back to \(C\) is the identity, as required.
\end{proof}

\begin{proposition}[The pretriangulated hull and perfect objects]\label{res:G12}
There is a small full dg category \(\operatorname{pretr}(\mathcal A)\)
of finite-cell semi-free modules that contains the Yoneda image,
represents all finite-cell semi-free modules up to dg isomorphism,
and is closed under zero, shifts, finite sums, and cones of closed
degree-zero maps. Its homotopy category embeds as a triangulated
subcategory of \(D(\mathcal A)\). The perfect subcategory,
obtained by also taking retracts in \(D(\mathcal A)\), is exactly
the subcategory of compact objects. The Yoneda map
\(\mathcal A\to\operatorname{pretr}(\mathcal A)\) is a quasi-equivalence
precisely when all finite-cell objects are represented, up to
homotopy, by objects of \(\mathcal A\).
\end{proposition}
\begin{proof}
The cone of a map from \(P\) to \(Q\) has graded module
\(Q\oplus P[1]\). Order all cells of \(Q\) first and then those
of \(P[1]\) in their existing order of stages. The shifted
\(P\)-differential lowers these latter stages, and the cone's
off-diagonal term lies in \(Q\). This orders the cone as a
finite-cell module. Shifts preserve the decreasing property of
the cell differential, and finite sums admit an ordering by
the original stages. Conversely attaching one cell of degree
\(n\) with boundary \(z\) is the cone of
\(h_a[-n-1]\to P\) carrying the generator to \(z\), by the cone
formula. Finite-cell modules are therefore exactly the finite
constructions generated by the representables in this model,
up to choices of graded cell presentations. For smallness,
start with the set of written representables and zero. At
each finite stage adjoin the written shifts, finite sums and
cones of every closed degree-zero map among the present objects,
using the formulas of Result~\ref{res:G04}. These data range
over sets because each Hom is a set. The union over stages
is a set closed under the stated operations. Take the full
dg subcategory on this set.

All finite-cell modules are \(K\)-projective. By
Result~\ref{res:G08} their homotopy morphisms compute the
derived morphisms, and by the cone calculation the full image
is triangulated. Result~\ref{res:G11} gives the assertion about
perfect objects. A dg functor is a quasi-equivalence if it
induces quasi-isomorphisms on all morphism complexes and is
essentially surjective on \(H^0\). Yoneda is already an
isomorphism on the morphism complexes, so the condition is
exactly the stated representability up to homotopy. In this
case the triangulation transports to \(H^0(\mathcal A)\).
\end{proof}

\section{Derived tensor functors and Morita equivalence}

Let \(\mathcal A,\mathcal B\) be small dg categories over \(k\).
An \(\mathcal A\)-\(\mathcal B\) dg bimodule will be described by a
dg functor \(T:\mathcal A\to \mathrm{dgMod}\text{-}\mathcal B\).
In particular \(u:a\to b\) acts as \(T(u):T(a)\to T(b)\) of the
same degree as \(u\), in the Hom convention of Result~\ref{res:G03}.

\begin{proposition}[The signed bimodule adjunction]\label{res:G13}
For a right dg \(\mathcal A\)-module \(M\), \(M\otimes_\mathcal A T\)
is the quotient of \(\bigoplus_a M(a)\otimes_k T(a)\) by the
homogeneous balancing relations
\((m u)\otimes t=m\otimes T(u)t\).
For a right dg \(\mathcal B\)-module \(N\), put
\(G(N)(a)=\underline{\operatorname{Hom}}_\mathcal B(T(a),N)\)
with right \(\mathcal A\)-action by precomposition. There is a
natural isomorphism of complexes
\[
 \underline{\operatorname{Hom}}_\mathcal B(M\otimes_\mathcal A T,N)
    \cong \underline{\operatorname{Hom}}_\mathcal A(M,G(N)).
\]
If each \(T(a)\) is \(K\)-projective, \(G\) sends quasi-isomorphisms
to quasi-isomorphisms. Computing the tensor on a \(K\)-projective
replacement yields an exact left adjoint
\(L_T:D(\mathcal A)\to D(\mathcal B)\) to \(G\) on derived categories.
Moreover \(L_T(h_a)\cong T(a)\).
\end{proposition}
\begin{proof}
The differential on the direct sum is the tensor differential.
It preserves the balancing relation because its derivative splits
into the three relations for \(dm,u,t\), \(m,du,t\), and \(m,u,dt\),
with coefficients \(1,(-1)^{|m|}\), and \((-1)^{|m|+|u|}\).
This uses
\(d(T(u)t)=T(du)t+(-1)^{|u|}T(u)dt\), which says that \(T\) is dg.
The right \(\mathcal B\)-action on the second tensor factor
respects balancing, since each \(T(u)\) is graded right-linear.
The tensor differential obeys the right Leibniz rule for this
action. For the right \(\mathcal A\)-module \(G(N)\), write
\(g u=g\circ T(u)\). The Hom composition rule of
Result~\ref{res:G03} gives its Leibniz rule.

If \(F:M\otimes_\mathcal A T\to N\) has degree \(r\), send it to
the graded map \(m\mapsto (t\mapsto F(m\otimes t))\). Balancing
is precisely \(\mathcal A\)-linearity of the curried map, and its
values are \(\mathcal B\)-linear maps. Conversely these two
linearity conditions define the balanced uncurried map.
For the differential, evaluating the curried differential
on \(m,t\) gives
\[
 d_N F(m\otimes t)-(-1)^r F(dm\otimes t)
       -(-1)^{r+|m|}F(m\otimes dt),
\]
which is the currying of the uncurried Hom differential.
This proves the isomorphism and its compatibility with closed
maps. The unit and counit are \(m\mapsto[t\mapsto m\otimes t]\)
and \(g\otimes t\mapsto g(t)\). Their composites on evaluated
tensors and maps are identities.

Now suppose the \(T(a)\) are \(K\)-projective. The formula of
Result~\ref{res:G08} for applying Hom from such modules shows
that \(G\) takes acyclics to acyclics and quasi-isomorphisms to
quasi-isomorphisms. For \(K\)-projective \(P\), the adjunction
of complexes shows \(P\otimes_\mathcal A T\) is \(K\)-projective
because its Hom into an acyclic \(N\) is Hom from \(P\) into the
acyclic \(G(N)\). The tensor takes homotopy equivalences to
homotopy equivalences by tensoring the homotopy with the identity.
Thus compute \(L_T\) on \(K\)-projective representatives.
The comparison of derived Hom groups in Result~\ref{res:G08}
now gives the derived adjunction. Tensoring commutes with shifts
and cones by the differential formulas, so \(L_T\) is exact.
For \(G\), the cone and shift in the target of Hom are likewise
the cone and shift of Hom, hence \(G\) is exact on the derived
categories. Finally \(h_a\otimes_\mathcal A T\to T(a)\) is
given by \(u\otimes t\mapsto T(u)t\), with inverse
\(t\mapsto 1_a\otimes t\). Balancing proves both composites are
identities and the maps commute with differentials.
\end{proof}

\begin{theorem}[A derived Morita criterion]\label{res:G14}
Assume each \(T(a)\) is \(K\)-projective, as in
Result~\ref{res:G13}. The derived adjunction is an equivalence
if the following conditions hold.
\begin{enumerate}
\item Each \(T(a)\) is compact in \(D(\mathcal B)\).
\item For every \(a,b\), the map of complexes
\[
  \mathcal A(a,b)\longrightarrow
     \underline{\operatorname{Hom}}_\mathcal B(T(a),T(b))
\]
induced by \(T\) is a quasi-isomorphism.
\item The shifted objects \(T(a)[n]\) detect zero in \(D(\mathcal B)\).
\end{enumerate}
Conversely, if \(L_T\) is an equivalence, these conditions hold.
\end{theorem}
\begin{proof}
Compactness and Result~\ref{res:G08} show that the natural map
\(\bigoplus_i G(N_i)\to G(\bigoplus_i N_i)\) is a quasi-isomorphism.
Indeed its cohomology at each \(a,n\) is the coproduct comparison
for maps from \(T(a)[-n]\), which is compact since \(T(a)\) is.
The functor \(L_T\) commutes with coproducts by the adjunction
or the tensor formula. Both functors are exact by
Result~\ref{res:G13}. In particular the full subcategory where
the unit \(X\to G(L_T(X))\) is an isomorphism is closed under
shifts, cones, isomorphisms, and coproducts. Closure under cones
follows by applying the Hom exact sequences to the comparison
of triangles. The unit on \(h_b\), evaluated at \(a\), is
exactly the map in condition (2), by the formula for the unit and
the evaluation \(h_b\otimes_\mathcal A T\cong T(b)\).
Thus this subcategory is all of \(D(\mathcal A)\) by
Result~\ref{res:G10}.

For \(Y\in D(\mathcal B)\) apply \(G\) to the counit
\(L_T G(Y)\to Y\). The triangular identity identifies this as
the inverse of the unit on \(G(Y)\). Condition (3) implies \(G\)
detects zero by computing its cohomology at each \(a\) using
maps from the shifts of \(T(a)\). Since it is exact it detects
isomorphisms by applying the zero detection to the cone.
Thus the counit is an isomorphism. This proves equivalence.

Conversely \(L_T\) carries each \(h_a\) to \(T(a)\).
An exact equivalence preserves coproducts and the compact
generation and detection properties in Result~\ref{res:G10}.
Full faithfulness on all shifted representables identifies the
cohomology in every degree of the map in condition (2) with
an isomorphism of derived morphism groups. This identification
uses the Yoneda evaluation on the source and the
\(K\)-projectivity of \(T(a)\) on the target. Hence that map is
a quasi-isomorphism.
\end{proof}

\begin{corollary}[Invariance under quasi-equivalence]\label{res:G15}
A quasi-equivalence \(F:\mathcal A\to\mathcal B\) induces an
equivalence of derived categories by extension of scalars, with
restriction along \(F\) as inverse.
\end{corollary}
\begin{proof}
Use \(T(a)=h_{F(a)}\) with its Yoneda functoriality. Enriched Yoneda
identifies \(G\) with restriction. Each \(T(a)\) is a compact
\(K\)-projective module. The quasi-isomorphism condition on
morphism complexes of \(F\) is condition (2) of
Result~\ref{res:G14}. Every \(b\) is isomorphic to some \(F(a)\)
in \(H^0(\mathcal B)\), so every \(h_b\) is isomorphic to
some \(T(a)\) in the derived category. The detection assertion
of Result~\ref{res:G10} gives condition (3).
\end{proof}

\begin{corollary}[An endomorphism dg algebra of a generator]\label{res:G16}
Let \(P\) be a compact \(K\)-projective right \(\mathcal B\)-module
whose shifts detect zero. Then \(D(\mathcal B)\) is equivalent
to \(D(\underline{\operatorname{End}}_\mathcal B(P))\).
Also the inclusion of \(\mathcal A\) into its pretriangulated
hull induces an equivalence of derived categories.
\end{corollary}
\begin{proof}
Regard \(\underline{\operatorname{End}}_\mathcal B(P)\) as a
one-object dg category acting on \(P\). Its map in condition (2)
of Result~\ref{res:G14} is the identity. The remaining conditions
are the hypotheses on \(P\), proving the first assertion.
For the hull, use Yoneda in \(\operatorname{pretr}(\mathcal A)\)
on each object \(h_a\) in the image of \(\mathcal A\).
The induced maps of morphism complexes are isomorphisms.
A finite-cell object of the hull is formed by shifts and cones
from the \(h_a\). The representable module of that finite-cell
object belongs to the triangulated subcategory generated by
the representables of \(h_a\). In fact Hom into a shift is a
shift of Hom, and Hom into a cone is the corresponding cone
objectwise, by the differential of Result~\ref{res:G03}.
Thus the representables of \(h_a\) detect zero in
\(D(\operatorname{pretr}(\mathcal A))\) by Result~\ref{res:G10}
and exact Hom from triangles. Apply Result~\ref{res:G14} again.
\end{proof}

For a concrete application take a field \(k\) and \(P=k^m\)
over \(k\) with \(m\geq2\). It is \(K\)-projective, compact,
and its shifts detect zero since \(k\) is a direct summand.
Its endomorphism dg algebra is \(\operatorname{Mat}_m(k)\)
in degree zero. Thus \(D(k)\) and \(D(\operatorname{Mat}_m(k))\)
are equivalent by Result~\ref{res:G16}. Any quasi-equivalence
of the corresponding one-object dg categories would give an
isomorphism of their degree-zero endomorphism \(k\)-algebras,
which have different dimensions. Even a zigzag of
quasi-equivalences induces equivalences of the \(H^0\) categories
and would give the same impossible isomorphism. This illustrates
the distinction between quasi-equivalence and derived Morita
equivalence. The criterion above concerns an explicit bimodule
functor, so its conclusion also records the dg structure used
to induce the equivalence.

The dg-category localization at Morita equivalences and the
comparison of arbitrary enhancements involve further structure.
See \cite[\S\S3.8, 4.5]{keller} for these developments.
